\section{Цифровой звук: I²S и микрофон}\label{les:i2s}
\lessonintro
  {получать звуковые отсчёты и измерять уровень сигнала}
  {уроки 6, 8 и 24}
  {график относительной громкости речи}
  {INMP441, провода; для расширения --- I²S-усилитель}

\subsection{Звук --- поток чисел}
Пищалка из урока 24 воспроизводила одну частоту. Для речи нужен поток отсчётов: каждое число
задаёт амплитуду в определённый момент. Частота дискретизации 16000 Гц означает 16000 отсчётов
в секунду. Разрядность задаёт точность отсчёта, число каналов --- количество потоков.
I²S передаёт такой поток между чипом и цифровым микрофоном, кодеком или усилителем.

I²S и I²C --- разные интерфейсы. У стандартного I²S есть BCLK (такты битов), WS (выбор канала)
и линия данных. В нашем опыте ESP32-S3 формирует такты, микрофон выдаёт данные.
Согласуйте частоту, длину слота, канал и формат Philips I²S. Для 16000 кадров/с, двух слотов
по 32 бита получаем BCLK $16000\cdot2\cdot32=1024000$ Гц.

\subsection{Подключение и первый поток}
Отключите предыдущую схему. INMP441: VDD к 3V3, GND к GND, SCK к GPIO16,
WS к GPIO17, SD к GPIO18, L/R к GND. Последняя перемычка выбирает левый слот.
Проверьте питание и маркировку конкретного модуля.

\begin{figure}[H]
\centering
\begin{circuitikz}[scale=0.95,transform shape]
  \node[blk,fill=blkcpu,minimum width=2.8cm,minimum height=4.6cm] (m) at (0,0) {};
  \node[blk,fill=blkper,minimum width=2.8cm,minimum height=4.6cm] (d) at (8,0) {};
  \node[font=\sffamily\small\bfseries] at (0,2.8) {ESP32-S3};
  \node[font=\sffamily\small\bfseries,align=center] at (8,2.95) {INMP441\\[-2pt]\scriptsize микрофонный модуль};
  \foreach \y/\n in {1.8/{\pin{16}},0.9/{\pin{17}},0/{\pin{18}},-0.9/3V3,-1.8/GND}
    \node[font=\sffamily\scriptsize,anchor=east] at (1.4,\y) {\n};
  \foreach \y/\n in {1.8/SCK,0.9/WS,0/SD,-0.9/VDD,-1.35/{L/R},-1.8/GND}
    \node[font=\sffamily\scriptsize,anchor=west] at (6.6,\y) {\n};
  \draw[thick,accent]   (1.4,1.8)  -- (6.6,1.8);
  \draw[thick,okgreen]  (1.4,0.9)  -- (6.6,0.9);
  \draw[thick,warn]     (1.4,0)    -- (6.6,0);
  \draw[thick,esp]      (1.4,-0.9) -- (6.6,-0.9);
  \draw[thick,black!70] (1.4,-1.8) -- (6.6,-1.8);
  \draw[thick,black!70] (6.6,-1.35) -- (5.6,-1.35) -- (5.6,-1.8);
  \fill[black!70] (5.6,-1.8) circle (1.5pt);
  \node[font=\sffamily\scriptsize,text=accent,fill=white,inner sep=1pt] at (4,1.8) {BCLK};
  \node[font=\sffamily\scriptsize,text=okgreen,fill=white,inner sep=1pt] at (4,0.9) {выбор канала};
  \node[font=\sffamily\scriptsize,text=warn,fill=white,inner sep=1pt] at (4,0) {данные};
\end{circuitikz}
\caption{Подключение микрофона INMP441: три сигнальные линии, питание \SI{3.3}{V} и перемычка L/R на землю
(левый слот). Порядок выводов на модуле зависит от производителя}
\end{figure}

В Arduino-ESP32 3.2.0 используем
\code{ESP\_I2S.h}, а не примеры старого драйвера \code{driver/i2s.h}.

\codecap{Отсчёты INMP441 и среднеквадратичный уровень}
\codeinput{l29_i2s_meter}

\subsection{Что именно измеряет пример}
У микрофона полезные данные занимают 24 бита внутри 32-битного слота. Деление на 256 убирает
младшие биты, сохраняя знак. Вычитаем среднее окна, чтобы постоянное смещение не считалось звуком.
Вычисляем $\mathrm{RMS}=\sqrt{\frac1N\sum(x_i-\bar x)^2}$.
Это относительный уровень, а не дБ SPL: для звукового давления нужна акустическая калибровка.

DMA передаёт блоки между периферией и памятью: CPU не обслуживает каждый бит.
Чтение блока может ждать данные. Для устройства с экраном и кнопками вынесите чтение аудио
в отдельную задачу, а результаты передавайте очередью. Буфер должен оставаться доступным,
пока им пользуется драйвер. Не печатайте каждый отсчёт текстом на 115200 бод: каналу может
не хватить пропускной способности.

\subsection{От измерения к воспроизведению}
Сравните фон, речь и хлопок на одинаковом расстоянии. Добавьте гистерезис индикатора «громко».
Для воспроизведения используют I²S-усилитель, например MAX98357A: к нему идут BCLK, WS и DOUT,
динамик подключается к выходам усилителя. Питание выбирают по документации модуля.
У ESP32-S3 нет встроенного аналогового DAC: соответствующие примеры первого ESP32 нельзя
переносить напрямую. Цифровой I²S, внешний DAC и PWM решают разные задачи.

\begin{trouble}
Нули: выбран не тот слот или отсутствуют такты. Шум без звука: формат, выравнивание или питание.
Пропуски: обработка и печать занимают слишком много времени. Проверяйте фактическое число
прочитанных байтов, а не предполагаемую длину массива.
\end{trouble}
\begin{summary}
Согласуем такты, слоты и формат, обрабатываем блоки. Различаем относительную амплитуду,
цифровой уровень и откалиброванное звуковое давление.
\end{summary}

\begin{selfcheck}
\begin{enumerate}
\item Чем I²S отличается от I²C?
\item Почему 24-битный микрофон читаем в 32-битный массив?
\item Почему RMS нельзя подписать «дБ SPL»?
\end{enumerate}
\end{selfcheck}
\needspace{20\baselineskip}
\begin{tasks}
\begin{enumerate}
\item \lv{1} Составьте таблицу соединений BCLK, WS, SD и питания.
\item \lv{1} Запишите фон и максимум хлопка на расстоянии 30 см.
\item \lv{2} Сравните окна 256 и 512 отсчётов по задержке обновления.
\item \lv{2} Добавьте индикатор «громко» с разными порогами включения и выключения.
\item \lv{3} Вынесите чтение в задачу и измерьте число пропущенных блоков при работе экрана.
\item \lv{4} Запишите PCM/WAV на SD: рассчитайте размер и проверьте длительность на компьютере.
\end{enumerate}
\end{tasks}

\paragraph{Документация для продолжения}\href{https://docs.espressif.com/projects/arduino-esp32/en/latest/api/i2s.html}{Официальное руководство по теме}.
